\scanpage{183}\mspage{38}%
% archive sheet number: 183
% Reader A. Main manuscript baselines L1--L23; no paragraph break securely established at L9.
dans cette première approche pré-catégorie dérivée, par la connaissance du groupe fondamental $\pi_1$, donc des catégories des coeff. \unsure{constant|constants} tordus $= \pi_1$-\unsure{module|modules}, les foncteurs
% L5: the separate point near (344,445), above and slightly left of the comma, has uncertain origin/attachment; it is not assigned a mathematical role. A semicolon association remains possible.
$H^i(-\raisebox{0pt}[0pt][0pt]{\parbox{0pt}{\figure{\begin{tikzpicture}[overlay]\fill (-0.03,0.23) circle (0.35pt);\end{tikzpicture}}}},M)$ sur ceux-ci avec \unsure{le|la} structure de \unsure{foncteur|foncteurs} \deleted{\unsure{ab|ab.}, de \unsure{foncteur|foncteurs}} \inserted{\unsure{cohomologique|cohomologiques}} relatif aux suites exactes, \unsure{le|la}
% L7: the older struc-like layer is superposed on the beginning of live structure; exact stroke chronology is insecure.
\deleted{\unsure{struc}}structure de cup-produit etc --%
% L8: detached compact point near (247,793), below the left end of the dash; origin and function uncertain.
\raisebox{0pt}[0pt][0pt]{\parbox{0pt}{\figure{\begin{tikzpicture}[overlay]\fill (-0.21,-0.12) circle (0.3pt);\end{tikzpicture}}}}
reliés par certaines propriétés formelles. Une fois qu'on dispose du langage des cat. dérivées, la \inserted{sous catégorie de la catégorie dérivée des} complexes abéliens, \unsure{forme|formée|formés|form-} \inserted{aux complexes de} faisceaux de coh. loc constante sur $X$, avec sa structure \unsure{triangulé|triangulée} et sa structure
% L13: a retraced/cancelled initial fragment precedes the live reading multiplicative; the fragment and its layering remain doubtful.
\deleted{\unsure{m|sou|de}}multiplicative
% L13: the distinct point near (1164,1373) lies well below the tensor sign, with no assigned mathematical meaning.
$\overset{\mathbb{L}}{\otimes}\raisebox{0pt}[0pt][0pt]{\parbox{0pt}{\figure{\begin{tikzpicture}[overlay]\fill (-0.12,-0.30) circle (0.4pt);\end{tikzpicture}}}}$ (et éventuellement $\mathbb{R}\mathrm{Hom}$ \ldots) donnant un candidat plus satisfaisant pour espérer récupérer le type d'homotopie. Je n'ai guère d'idée
% L17: a short detached dark stroke near (241,1835) lies below si; origin and function uncertain.
si\raisebox{0pt}[0pt][0pt]{\parbox{0pt}{\figure{\begin{tikzpicture}[overlay]\draw[line width=0.7pt] (-0.08,-0.26) -- (-0.065,-0.29) -- (-0.075,-0.32);\end{tikzpicture}}}} oui ou non ils permettent de le récupérer \inserted{en effet}, mais par contre%
% L18: two separate compact points near (943,1958) and (938,1966), below/left of contre; no punctuation function or hand is assigned.
\raisebox{0pt}[0pt][0pt]{\parbox{0pt}{\figure{\begin{tikzpicture}[overlay]\fill (-0.55,-0.30) circle (0.3pt);\fill (-0.59,-0.36) circle (0.3pt);\end{tikzpicture}}}}
n'ai guère de doute qu'en allant au bout de la ``linéarisation'', c'est à dire en partant dans le
% L21: the single underline is physically in two segments, approximately (317,2274)--(557,2261) and (582,2262)--(671,2265). It begins within the first n and resumes within the horizontal extent of l. The macro records the covered word span; the gap is not interpreted as exclusion of l from emphasis.
\uline{non abélien} et travaillant avec la $(n+1)$-catégorie (\unsure{sans|sa} \unsure{structures|structure} \unsure{supplémentaires|supplémentaire} dessus!) des $n$-champs loc. constants
% L23a: dark raised sur X has a thin curved guide to the interval after constants and before pour. Reading, retracing history, and possible cancellation remain qualified; no distinct cancelled word is recoverable.
\inserted{\unsure{sur $X$}} pour tout $n$
